\section{Опис вимог}

\subsection{Сценарій місії}

Операційний сценарій системи \projname{}:

\begin{enumerate}[itemsep=3pt]
  \item Оператор запускає скрипт місії (\texttt{missions/relay\_demo.py} або \texttt{missions/inspection.py}) з наземної станції.
  \item Усі три дрони (\texttt{uav-01}, \texttt{uav-02}, \texttt{uav-03}) злітають до висоти $z=3{,}0$~м — на 1~м вище верхівок стелажів.
  \item Кожен дрон летить до точки інтересу (POI) свого ряду стелажів, зависає на 5~с (імітація детального знімання) та виконує орбіту радіусом 1,5~м навколо POI.
  \item Дрон виконує прямолінійний прохід (Boustrophedon sweep) вздовж осі~Y свого коридора від $y=-5{,}5$~м до $y=+5{,}5$~м.
  \item При сценарії естафети (\texttt{relay\_demo.py}): коли заряд батареї \texttt{uav-01} падає нижче порогу (за замовчуванням 60\%), \texttt{uav-02} перехоплює поточну позицію і продовжує сканування.
  \item Після завершення всіх проходів дрони сідають на початкові позиції.
  \item Grafana відображає прогрес в режимі реального часу.
\end{enumerate}

\subsection{Функціональні вимоги}

\begin{enumerate}[label=\textbf{ФВ-\arabic*}., itemsep=4pt]
  \item Система повинна підтримувати одночасний незалежний політ не менше ніж 3 БПЛА.
  \item Кожен дрон повинен реагувати на команди оператора: \texttt{takeoff}, \texttt{goto}, \texttt{orbit}, \texttt{inspect\_poi}, \texttt{scan}, \texttt{hover}/\texttt{stop}, \texttt{land}.
  \item Команда \texttt{stop} повинна перервати поточну траєкторію і зафіксувати дрон у поточній позиції не пізніш ніж за 0{,}5~с після отримання.
  \item Система повинна публікувати телеметрію кожного дрона з частотою $\geq 5$~Гц у топік \texttt{swarm/telemetry/\{drone\_id\}}.
  \item При падінні заряду батареї нижче порогу система \emph{автоматично} (без втручання оператора) ініціює процедуру естафетної передачі місії.
  \item Дрон-наступник повинен розпочати сканування не далі ніж на 0{,}5~м від позиції зупинки попередника.
  \item Grafana дашборд повинен відображати: поточний режим, заряд батареї у \%, прогрес покриття у \% для кожного дрона.
\end{enumerate}

\subsection{Нефункціональні вимоги}

\begin{enumerate}[label=\textbf{НВ-\arabic*}., itemsep=4pt]
  \item \textbf{Науковість}: усі алгоритми (керування, траєкторія, покриття) мають бути обґрунтовані посиланнями на рецензовані публікації.
  \item \textbf{Точність}: відхилення фактичної траєкторії від референсної по осі~X (поперечне) не повинно перевищувати 0{,}1~м при прямолінійному проході.
  \item \textbf{Тестованість}: ключові модулі повинні мати одиничні тести (coverage $\geq 80$\% по модулях логіки).
  \item \textbf{Переносимість}: система повинна запускатися як у симуляторі Isaac Sim + Pegasus, так і у headless-режимі (\texttt{mock\_swarm}) без зміни коду місії.
  \item \textbf{Конфігурованість}: ключові параметри (швидкість, висота, поріг батареї, drain rate) повинні задаватися через файл \texttt{.env} без зміни коду.
\end{enumerate}

\subsection{Карта об'єкта інспекції}

Координатна система Isaac Sim (рис.~\ref{fig:facility_map}): X — вправо, Y — вглибину, Z — вгору; підлога $Z=0$.

\begin{table}[h]
  \centering
  \caption{Точки інтересу (POI) та зони сканування}
  \begin{tabularx}{\linewidth}{lllX}
    \toprule
    Ідентифікатор & Тип & Координати (X, Y, Z), м & Призначений дрон \\
    \midrule
    \texttt{rack\_left\_end}   & POI      & ($-3{,}25$; $+5{,}5$; $3{,}0$)  & \texttt{uav-01} \\
    \texttt{rack\_center\_end} & POI      & ($0{,}00$; $+5{,}5$; $3{,}0$)   & — \\
    \texttt{rack\_right\_end}  & POI      & ($+3{,}25$; $+5{,}5$; $3{,}0$)  & \texttt{uav-02} \\
    \texttt{aisle\_left}   & ScanZone & $X=-3{,}25$, $Y\in[-5{,}5;+5{,}5]$ & \texttt{uav-01} \\
    \texttt{aisle\_center} & ScanZone & $X=+3{,}25$, $Y\in[-5{,}5;+5{,}5]$ & \texttt{uav-02} \\
    \texttt{aisle\_outer}  & ScanZone & $X=+5{,}75$, $Y\in[-5{,}5;+5{,}5]$ & \texttt{uav-03} \\
    \bottomrule
  \end{tabularx}
  \label{tab:facility}
\end{table}
